\section*{Chapitre \ref{ch:g12:diffeq} --- Équations différentielles}

\begin{solution}{exo:g12:diffeq:1}
\emph{(a)} $y = C\eu^{3x}$~; $y(0) = 2$ donne $y = 2\eu^{3x}$.

\emph{(b)} $y' = -\frac12 y$, donc $y = C\eu^{-x/2}$~; $y(0) = -1$ donne
$y = -\eu^{-x/2}$.

\emph{(c)} Équilibre $y_p = 5$~; $y = C\eu^{-x} + 5$~; $y(0) = 0$ donne
$C = -5$~: $y = 5\left(1 - \eu^{-x}\right)$.
\end{solution}

\begin{solution}{exo:g12:diffeq:2}
$N' = kN$, donc $N(t) = N_0 \eu^{kt}$. Doubler en $3$~heures signifie
$\eu^{3k} = 2$, \emph{c.-à-d.}
\[
k = \frac{\ln 2}{3} \approx 0.231\ \text{h}^{-1}.
\]
\end{solution}

\begin{solution}{exo:g12:diffeq:3}
$y_p'(x) = \eu^x + x\eu^x = y_p(x) + \eu^x$~: $y_p$ est une solution
particulière (c'est le cas résonant de la
\cref{met:g12:diffeq:particular}). Par la \cref{prop:g12:diffeq:general},
les solutions sont $y = C\eu^{x} + x\eu^{x} = (C + x)\,\eu^x$,
$C \in \R$.
\end{solution}

\begin{solution}{exo:g12:diffeq:4}
Essayer $y_p = \beta\eu^x$~: $\beta\eu^x = -2\beta\eu^x + \eu^x$ donne
$3\beta = 1$, donc $y_p = \frac13\eu^x$. Solution générale
$y = C\eu^{-2x} + \frac13\eu^{x}$~; la condition $y(0) = 1$ donne
$C = \frac23$~:
\[
y(x) = \frac{2}{3}\,\eu^{-2x} + \frac{1}{3}\,\eu^{x}.
\]
\end{solution}

\begin{solution}{exo:g12:diffeq:5}
$N(t) = N_0\,\eu^{-\lambda t}$ avec $\lambda = \frac{\ln 2}{5730}$
(\cref{ex:g12:exp:halflife}). On résout $\eu^{-\lambda t} = 0.6$~:
\[
t = \frac{\ln(1/0.6)}{\lambda} = 5730\,\frac{\ln(5/3)}{\ln 2}
\approx 5730 \times \frac{0.5108}{0.6931} \approx 4220 \text{ ans}.
\]
\end{solution}

\begin{solution}{exo:g12:diffeq:6}
\emph{1.} Le sel entre à $0.2 \times 5 = 1$~kg/min. Le flux sortant emporte
la concentration $\frac{m}{100}$~kg/L à $5$~L/min, \emph{c.-à-d.}
$\frac{m}{20}$~kg/min. D'où $m' = 1 - \frac{m}{20}$.

\emph{2.} Équilibre $m_p = 20$~; $m(t) = 20 + C\eu^{-t/20}$, et
$m(0) = 0$ donne $C = -20$~:
\[
m(t) = 20\left(1 - \eu^{-t/20}\right) \xrightarrow[t\to+\infty]{} 20 \text{ kg}.
\]
À long terme la concentration du réservoir égale celle de la saumure
entrante~: $0.2$~kg/L $\times$ $100$~L $= 20$~kg.
\end{solution}

\begin{solution}{exo:g12:diffeq:7}
\emph{1.} $v' = g - \frac{k}{m} v$ avec $\frac km = \frac{16}{80} = 0.2$.
Équilibre $v_\infty = \frac{mg}{k} = \frac{9.8}{0.2} = 49$~m/s~;
$v(t) = 49 + C\eu^{-0.2t}$, et $v(0) = 0$ donne
\[
v(t) = 49\left(1 - \eu^{-0.2 t}\right).
\]

\emph{2.} Vitesse limite $49$~m/s ($\approx 176$~km/h). On veut
$1 - \eu^{-0.2t} = 0.95$, \emph{c.-à-d.} $\eu^{-0.2t} = 0.05$~:
\[
t = \frac{\ln 20}{0.2} \approx \frac{3.00}{0.2} \approx 15 \text{ s}.
\]
\end{solution}

\begin{solution}{exo:g12:diffeq:8}
\emph{1.} $z = \frac1y$ donne $z' = -\frac{y'}{y^2}
= -\frac{y(1-y)}{y^2} = -\frac{1-y}{y} = -\frac1y + 1 = -z + 1$.

\emph{2.} Équilibre $z_p = 1$, donc $z(t) = 1 + C\eu^{-t}$. De
$y(0) = \frac{1}{10}$, $z(0) = 10$, donc $C = 9$ et
\[
y(t) = \frac{1}{1 + 9\,\eu^{-t}} .
\]

\emph{3.} Lorsque $t \to +\infty$, $9\eu^{-t} \to 0$ et $y(t) \to 1$. La
courbe est la classique courbe logistique en S (\emph{sigmoïde})~:
croissance lente au début ($y$ petit, $y' \approx y$), croissance la plus
rapide lorsque $y = \frac12$ (où $y' = y(1-y)$ est maximale), puis
saturation vers la capacité $1$.
\end{solution}

\begin{solution}{exo:g12:diffeq:9}
Par le \cref{thm:g12:diffeq:affine}, $y(x) = C\eu^{ax} - \frac ba$. D'où
\[
u_{n+1} = C\eu^{a(n+1)} - \frac ba
= \eu^{a}\left(C\eu^{an} - \frac ba\right) + \frac ba\left(\eu^a - 1\right)
= q\,u_n + r
\]
avec $q = \eu^{a}$ et $r = \frac{b}{a}\left(\eu^{a} - 1\right)$.
La condition $\abs q < 1$ signifie $\eu^a < 1$, \emph{c.-à-d.} $a < 0$~:
exactement la condition sous laquelle les solutions de l'\omterm{def:g12:diffeq:ode}{équation
différentielle} convergent vers l'équilibre $-\frac ba$ --- et en effet le
\omterm{pb:g10:functions:1}{point fixe} de la récurrence est $\frac{r}{1 - q} = -\frac{b}{a}$.
\end{solution}

\begin{solution}{pb:g12:diffeq:1}
\textbf{1.} $y = 2\eu^{3t}$. Pour la seconde~: équilibre $3$, donc
$y = 3 + C\eu^{-2t}$ avec $y(0) = 0$, d'où $C = -3$ et
$y = 3\left(1 - \eu^{-2t}\right)$.

\textbf{2.} $\left(y\eu^{-at}\right)' = y'\eu^{-at} -
ay\eu^{-at} = (y' - ay)\eu^{-at} = 0$~: le produit est une
constante $C$, donc $y = C\eu^{at}$ --- aucune solution
n'échappe.

\textbf{3.} Équilibre~: $y \equiv 3$. Toute autre solution s'écrit
$3 + C\eu^{-2t}$ avec $C \neq 0$~: l'exponentielle meurt et la
solution glisse vers $3$ --- un \omterm{pb:g10:functions:1}{point fixe} stable, jumeau continu
des \omterm{pb:g10:functions:1}{points fixes} de récurrence du \cref{pb:g11:seq:1}.

\textbf{4.} $y = y_0\eu^{-t/\tau}$ est divisé par deux lorsque
$\eu^{-t/\tau} = \frac12$, c'est-à-dire $t = \tau\ln 2$.

\textbf{5.} Toutes les solutions sont des translatées verticales de
la décroissance vers $3$~: celles qui partent au-dessus retombent
vers $3$, celles qui partent au-dessous montent vers $3$, et la
solution constante $3$ ne bouge pas --- un entonnoir autour de
l'équilibre.

\textbf{6.} $\lambda = \frac{\ln 2}{5730} \approx 1.21 \times
10^{-4}$ par an.

\textbf{7.} $\eu^{-\lambda t} = 0.2$~:
$t = \frac{\ln 5}{\lambda} \approx 13\,300$ ans.

\textbf{8.} $t = \frac{\ln(1/0.15)}{\lambda} \approx 15\,700$
ans~: les taureaux de Lascaux datent de la fin de la période
glaciaire.

\textbf{9.} $t = \frac{\ln(1/0.53)}{\lambda} \approx 5\,250$ ans~:
à une vie d'homme près de la valeur des archéologues --- l'horloge
fonctionne.

\textbf{10.} Dix demi-vies laissent $2^{-10} \approx 0.1\,\%$~: à
la limite de la mesure. Après $65$ millions d'années, la fraction
vaut $2^{-65\,000\,000/5\,730} \approx 2^{-11\,344}$~: il ne reste
aucun atome du stock initial dans le moindre fossile --- les
dinosaures se datent avec des horloges plus lentes
(potassium-argon et apparentées).

\textbf{11.} $52\,\%$ donne environ $5\,405$ ans, $54\,\%$ environ
$5\,095$~: le $\pm 1\,\%$ de la chimie devient à peu près
$\pm 155$ ans d'histoire --- la précision au laboratoire est la
précision du cartel de musée.

\textbf{12.} $z = T - T_a$ vérifie $z' = -kz$, donc
$z = (T_0 - T_a)\eu^{-kt}$ et $T = T_a + z$.

\textbf{13.} L'écart initial vaut $T_0 - T_a = 70$~; après cinq
minutes il vaut $70 - 20 = 50$, donc $50 = 70\eu^{-5k}$ et
$k = \frac{\ln(70/50)}{5} \approx 0.067$ par minute. Buvable~:
$55 - 20 = 35 = 70\eu^{-kt}$, d'où
$t = \frac{\ln 2}{k} \approx 10.3$ minutes.

\textbf{14.} La vitesse de perte est proportionnelle à l'écart
$T - T_a$~: un café brûlant perd sa chaleur le plus vite. Verser le
lait \emph{tout de \omterm{def:g12:seq:sequence}{suite}} réduit immédiatement l'écart, donc moins
de chaleur est perdue sur les dix minutes~; le verser à la fin
laisse d'abord le café refroidir à pleine vitesse. Pour la tasse la
plus chaude~: le lait d'abord. (Les hôtes impatients ont la
thermodynamique à l'envers.)

\textbf{15.} Écarts au \omterm{prop:g10:coordgeom:midpoint}{milieu} ambiant~: $10$ à minuit, $8$ une
heure plus tard, donc $\eu^{-k} = 0.8$ et
$k = \ln\frac{10}{8} \approx 0.223$ par heure. À la mort, l'écart
valait $17$~; il a décru jusqu'à $10$ à minuit~:
$s = \frac{\ln(17/10)}{k} \approx 2.4$ heures. Heure du décès~:
vers $21{:}40$.

\textbf{16.} Une pièce chauffée ou traversée de courants d'air
($T_a$ non constant) fausse l'horloge dans un sens ou dans
l'autre~; des vêtements ou une corpulence différente changent $k$
(les tables des légistes le corrigent) --- un $k$ erroné dilate tout
l'\omterm{def:g10:numbers:interval}{intervalle}~; et un corps déplacé depuis un autre lieu remet $T_a$
à zéro en cours de refroidissement, simulant une mort plus précoce
ou plus tardive. L'\omterm{def:g10:algebra:equation}{équation} est honnête~; ce sont les hypothèses qui
portent le risque.

\textbf{17.} $z = \frac1y$ obéit à $z' = 1 - z$~:
$z = 1 + 9\eu^{-t}$ (à partir de $z(0) = 10$), donc
$y = \frac{1}{1 + 9\eu^{-t}}$. Inflexion en $y = \frac12$~:
$9\eu^{-t} = 1$, soit $t = \ln 9 \approx 2.2$ --- la date de
retournement de l'épidémie, tirée directement de la formule.

\textbf{18.} Vitesse limite~: $v' = 0$ pour $v = 20$ m/s.
Solution~: $v(t) = 20\left(1 - \eu^{-t/2}\right)$. Puis $0.95$~:
$\eu^{-t/2} = 0.05$, donc $t = 2\ln 20 \approx 6$ s --- les gouttes
de pluie atteignent leur vitesse de croisière en quelques secondes,
et c'est pourquoi la pluie ne tue pas.

\textbf{19.} Régime permanent~: $c = 8$. La moitié~:
$4 = 8\left(1 - \eu^{-t/2}\right)$, donc $\eu^{-t/2} = \frac12$ et
$t = 2\ln 2 \approx 1.4$ heure. La perfusion est le jumeau continu
de l'emprunt~: entrée constante, sortie proportionnelle ---
l'\cref{exo:g12:diffeq:9} rend le dictionnaire exact.

\textbf{20.} Désintégration ($a < 0$, $b = 0$), refroidissement
($y = T - T_a$), chute ($a < 0$, $b > 0$~: approche de l'équilibre
par en dessous), perfusion (idem, dans une veine)~: une \omterm{def:g10:algebra:equation}{équation},
quatre mondes. La boucle~: énoncer la \omterm{def:g11:prob:rv}{loi} d'évolution~; écrire
l'\omterm{def:g10:algebra:equation}{équation}~; résoudre (\cref{met:g12:diffeq:solve})~; calibrer les
constantes sur des mesures~; prévoir~; puis auditer les hypothèses
(question 16). L'oscillation, elle, réclame $y'' = -\omega^2 y$ ---
rencontrée, résolue et mise en branle au \cref{pb:g12:trigo:1}.
\end{solution}
